% SIAM 期刊模板 —— SIAM J. Math. Anal. / SIAM J. Numer. Anal. / SIMAX
% 引擎：pdflatex
% 官方类 siamart 已安装时自动使用正式期刊格式（tlmgr install siam）；
% 未安装时降级为 article 预览排版，正文首段会给出提示，仍可正常编译与双向同步。
\IfFileExists{siamart.cls}{%
  \documentclass[siam,11pt]{siamart}
}{%
  \documentclass[11pt,a4paper]{article}
  \def\YTfallback{1}
}

\usepackage{amsmath,amssymb,amsthm}
\usepackage{graphicx}

\newtheorem{theorem}{Theorem}[section]
\newtheorem{lemma}[theorem]{Lemma}
\newcommand{\R}{\mathbb{R}}
\newcommand{\Om}{\Omega}
\newcommand{\eps}{\varepsilon}
\newcommand{\dd}{\,\mathrm{d}}

\title{Convergence of a Finite Element Scheme for a Nonlocal Allen--Cahn Equation}
\author{Yang Jianwei\thanks{School of Mathematics and Statistics, NCWU, Zhengzhou, P.R. China.}}
\date{}

\begin{document}
\maketitle

\ifdefined\YTfallback
\noindent\fbox{\parbox{0.96\linewidth}{\small\textbf{[YangTex preview mode] }
The official \texttt{siamart.cls} is not installed on this machine, so this
document is typeset with a generic article class. Install it with
\texttt{tlmgr install siam} (or the MiKTeX Console) and recompile: YangTex
will then switch to the real SIAM journal layout automatically.}}
\bigskip
\fi

\begin{abstract}\label{abs:siam}
We propose and analyse a fully discrete finite element scheme for a
nonlocal Allen--Cahn equation. Optimal-order error estimates in
$L^\infty(0,T;L^2)$ are proved under a mild time-step restriction.
\end{abstract}

\noindent\textbf{AMS subject classifications:} 65M60, 65M12, 35K55.

\section{Introduction}
Let $\Om\subset\R^d$, $d\le3$, be a convex polygonal domain. We consider
\begin{equation}\label{eq:ac}
  \partial_t u + \mathcal{L}u + \frac{1}{\eps^2}(u^3-u) = 0,
  \qquad (x,t)\in\Om\times(0,T],
\end{equation}
where $\mathcal{L}$ is a nonlocal operator of fractional type.
For the classical background and related results, see~\cite{AC}.

\begin{theorem}
Let $u$ and $U_h^n$ be the exact and fully discrete solutions. If
$\tau\le C h^{2}$, then
\begin{equation}
  \max_{1\le n\le N}\|u(t_n)-U_h^n\|_{L^2(\Om)}
  \le C(\tau + h^{2}).
\end{equation}
\end{theorem}

\section{Numerical analysis}
The proof uses a discrete Gronwall argument together with the
stability of the $L^2$ projection.

\bibliographystyle{siamplain}
\begin{thebibliography}{9}
\bibitem{AC} S.~M.~Allen and J.~W.~Cahn, \emph{A microscopic theory for
antiphase boundary motion}, Acta Metall. \textbf{27} (1979), 1085--1095.
\end{thebibliography}

\end{document}
